% \newpage
% \section*{Appendix: Stage 1 Brainstorming Notes}
% \addcontentsline{toc}{section}{Appendix: Stage 1 Brainstorming Notes}

% During the initial project phase, our team of six conducted a structured ideation session to decompose the Generative AI dilemma. We utilized two primary frameworks: the 5W1H method for problem framing and SCAMPER for solution ideation. Below are the transcribed notes from our collaborative session.

% \subsection*{1. Problem Framing via 5W1H Matrix}
% \begin{table}[h]
%     \centering
%     \renewcommand{\arraystretch}{1.5}
%     \begin{tabular}{|p{2cm}|p{12cm}|}
%         \hline
%         \textbf{Dimension} & \textbf{Raw Brainstorming Ideas \& Keywords} \\
%         \hline
%         \textbf{Who} & AI Corporations (OpenAI, DeepMind), Data Engineers, Independent Artists, Content Creators, Vulnerable Citizens (Deepfake victims), Lawmakers. \\
%         \hline
%         \textbf{What} & The conflict between unregulated web scraping (stealing data for AI training), intellectual property theft, and the creation of hyper-realistic fake media. \\
%         \hline
%         \textbf{When} & Present day -- The "Wild West" era of AI, where technological capabilities have drastically outpaced international copyright laws and detection tools. \\
%         \hline
%         \textbf{Where} & The global digital ecosystem, social media platforms (X, TikTok), digital art marketplaces, and mainstream news outlets. \\
%         \hline
%         \textbf{Why} & AI models need massive data to improve accuracy (Tech KPI). Artists need to protect their livelihood (Ethics/Finance). Society needs verifiable truth (Social Stability). \\
%         \hline
%         \textbf{How} & How can we build an ethical bridge? \textit{Initial Ideas:} Blanket bans? (Too extreme). Open-source only? (Dangerous). Opt-in licensing? (Hard to scale). Watermarking? (Technically challenging). \textbf{Chosen path: Fair Trade AI + Universal Watermarking.} \\
%         \hline
%     \end{tabular}
%     \caption{5W1H Brainstorming Matrix for Generative AI Dilemma}
%     \label{tab:5w1h}
% \end{table}

% \subsection*{2. Solution Ideation via SCAMPER}
% After framing the problem, we applied the SCAMPER technique to brainstorm potential engineering and policy solutions:
% \begin{itemize}
%     \item \textbf{Substitute:} Can we substitute scraped human data with AI-generated synthetic data? \textit{(Critique: This leads to "Model Collapse" over generations. Not a sustainable solution.)}
%     \item \textbf{Combine:} Can we combine blockchain technology with AI data ingestion? \textit{(Idea: Create a smart-contract royalty ledger to automatically pay artists when their opt-in data is used.)}
%     \item \textbf{Adapt:} How can we adapt traditional copyright watermarks for the AI age? \textit{(Idea: Invisible, pixel-level cryptographic watermarks like Google's SynthID.)}
%     \item \textbf{Modify:} Can we modify hardware to enforce ethics? \textit{(Idea: The "Hardware-Level Attestation" concept. GPUs must automatically sign generated outputs.)}
%     \item \textbf{Eliminate:} What happens if we eliminate the use of copyrighted data entirely? \textit{(Result: Dataset Shrinkage. Models become less capable but 100\% ethically compliant. We accept this trade-off.)}
% \end{itemize}


















\newpage
\section*{Appendix: Stage 1 Brainstorming Notes}
\addcontentsline{toc}{section}{Appendix: Stage 1 Brainstorming Notes}

During the initial stage of the project, our six-member team conducted a structured ideation session to clarify the problem space surrounding Generative AI in digital content creation. Two brainstorming frameworks were used: the \textbf{5W1H method} for problem framing and the \textbf{SCAMPER technique} for early-stage solution exploration. The notes below summarize the main directions generated during the collaborative discussion.

\subsection*{1. Problem Framing via 5W1H Matrix}
\begin{table}[h]
    \centering
    \renewcommand{\arraystretch}{1.5}
    \begin{tabular}{|p{2.2cm}|p{11.6cm}|}
        \hline
        \textbf{Dimension} & \textbf{Brainstorming Notes and Key Ideas} \\
        \hline
        \textbf{Who} & AI companies and model developers, data engineers, independent artists, content creators, digital platforms, vulnerable individuals affected by deepfakes, and lawmakers/regulators. \\
        \hline
        \textbf{What} & A growing conflict between large-scale web data collection for AI training, intellectual property protection, and the generation of highly realistic synthetic media that can be beneficial or harmful depending on use. \\
        \hline
        \textbf{When} & The present phase of rapid GenAI commercialization, in which technical capabilities are evolving faster than legal frameworks, policy adaptation, and detection tools. \\
        \hline
        \textbf{Where} & The global digital ecosystem, especially social media platforms, digital art marketplaces, online content-sharing environments, and news-related information spaces. \\
        \hline
        \textbf{Why} & AI organizations require large datasets to improve model capability and maintain competitiveness; creators seek protection of authorship, compensation, and recognition; society requires trustworthy digital information and protection from misuse. \\
        \hline
        \textbf{How} & How can GenAI systems be developed and governed to balance innovation, creator rights, and public trust? Early directions considered included stricter licensing models, provenance tracking, watermarking, platform-level safeguards, and access control mechanisms. \\
        \hline
    \end{tabular}
    \caption{5W1H Brainstorming Matrix for the Generative AI Content-Creation Dilemma}
    \label{tab:5w1h}
\end{table}

\subsection*{2. Early Solution Ideation via SCAMPER}
After the problem was framed, the team used the SCAMPER technique to explore possible engineering and governance directions. These ideas were not treated as final solutions, but as exploratory options for later evaluation.

\begin{itemize}
    \item \textbf{Substitute:} Can large-scale scraped copyrighted data be partially substituted with licensed, public-domain, or responsibly curated datasets? \textit{Observation: This may improve legal and ethical compliance, but could reduce dataset diversity and model performance.}

    \item \textbf{Combine:} Can AI data pipelines be combined with royalty-tracking or smart-contract mechanisms? \textit{Idea: A compensation layer could be added for opt-in creators whose works are used in training or referenced in commercial ecosystems.}

    \item \textbf{Adapt:} Can traditional attribution or watermarking practices be adapted for synthetic media? \textit{Idea: Invisible or cryptographic watermarking could help improve provenance and traceability of AI-generated outputs.}

    \item \textbf{Modify:} Can the design of GenAI systems be modified to include stronger safety or provenance controls by default? \textit{Idea: Generation systems could embed output signatures or metadata to support accountability and downstream verification.}

    \item \textbf{Put to another use:} Can the same generative technologies also be used defensively? \textit{Idea: AI could support detection, moderation, or authenticity verification systems rather than only content generation.}

    \item \textbf{Eliminate:} What happens if copyrighted or uncertain-source data is eliminated entirely from training pipelines? \textit{Observation: This may increase compliance and reduce disputes, but could also shrink training data availability and affect model capability.}

    \item \textbf{Reverse / Rearrange:} Instead of assuming unrestricted model access first and adding safeguards later, can governance be reversed so that provenance, consent, and accountability are built into the pipeline from the beginning? \textit{Idea: A “compliance-by-design” approach may be more sustainable than retroactive regulation.}
\end{itemize}

\subsection*{3. Brainstorming Outcome}
The brainstorming session helped the team move from a broad discussion about ``AI versus human creators'' toward a more structured understanding of the issue as a \textbf{socio-technical engineering challenge}. In particular, the 5W1H analysis clarified the affected stakeholders, the urgency of the problem, and the competing interests involved, while SCAMPER helped identify possible intervention points in data governance, content provenance, platform responsibility, and system design. These notes informed the more formal \textit{Brainstorming \& Problem Framing} section in the main report.